% Generated draft: requires mathematical and source review.
\documentclass[11pt]{article}
\usepackage[margin=25mm]{geometry}
\usepackage{fontspec}
\setmainfont{Latin Modern Roman}
\usepackage{amsmath,amssymb,mathtools,unicode-math}
\setmathfont{Latin Modern Math}
\usepackage{enumitem,needspace,xurl,hyperref}
\setlength{\parindent}{0pt}
\setlength{\parskip}{5pt}
\setlength{\emergencystretch}{3em}
\newcommand{\N}{\mathbb{N}}
\newcommand{\Z}{\mathbb{Z}}
\newcommand{\Q}{\mathbb{Q}}
\newcommand{\R}{\mathbb{R}}
\newcommand{\C}{\mathbb{C}}
\newcommand{\F}{\mathbb{F}}
\newcommand{\A}{\mathbb{A}}
\newcommand{\PP}{\mathbb{P}}
\newcommand{\E}{\mathbb{E}}
\newcommand{\eps}{\varepsilon}
\newcommand{\dd}{\,\mathrm{d}}
\begin{document}
\section*{UnsolvedMath ID: 30004309}
\textbf{Source problem code:} catalogue entry
\subsection{Definitions and mathematical statement}
A \emph{skew left brace} is a set \(B\) with group operations \(+\) and \(\circ\), called the additive and multiplicative operations, respectively, satisfying
\[
 a\circ(b+c)=(a\circ b)+(-a)+(a\circ c)
 \qquad(a,b,c\in B),
\]
where \(-a\) denotes the inverse in \((B,+)\); the additive group need not be abelian. The brace is \emph{two-sided} if it also satisfies
\[
 (a+b)\circ c=(a\circ c)+(-c)+(b\circ c)
 \qquad(a,b,c\in B).
\]
It is of \emph{abelian type} if \((B,+)\) is abelian. A finite group is an \emph{IYB group} (involutive Yang--Baxter group) if it is isomorphic to the multiplicative group of a finite skew left brace of abelian type; two-sidedness is not required in this definition.

For a group \(G\), write \([G,G]\) for the subgroup generated by the commutators \([x,y]=x^{-1}y^{-1}xy\). Define \(G^{(0)}=G\) and \(G^{(n+1)}=[G^{(n)},G^{(n)}]\). The group is \emph{solvable} if \(G^{(n)}=\{1\}\) for some integer \(n\geq0\). Nilpotency class exactly two means
\[
 \{1\}\neq [G,G]\subseteq Z(G),
\]
where \(Z(G)\) is the center. For a prime \(p\mid |G|\), a Sylow \(p\)-subgroup has order equal to the largest power of \(p\) dividing \(|G|\).

The supplied record asks the following three questions.
\begin{enumerate}
\item For every finite skew left brace \((B,+,\circ)\), does solvability of \((B,+)\) imply solvability of \((B,\circ)\)?
\item For every finite group \(G\) of nilpotency class two, does there exist a finite two-sided skew brace \((B,+,\circ)\) such that \((B,+)\) is abelian and \((B,\circ)\cong G\)?
\item Does there exist a finite group \(G\) which is not an IYB group but for which every Sylow \(p\)-subgroup is an IYB group, for every prime \(p\mid |G|\)?
\end{enumerate}
Here class two is read as exactly two. If the source intended class at most two, the additional abelian cases in question 2 admit the brace obtained by taking both operations to be the given abelian group operation, so this convention does not change the substantive question. The source's report-year discrepancy is recorded in Sources and remains unresolved. The supplied open-status and literature assessments are dated curation, not independent verification.
\subsection{Short English statement}
The questions ask whether solvability passes from the additive group of a finite skew left brace to its multiplicative group, whether every finite group of nilpotency class two admits a two-sided brace realization with abelian additive group, and whether a finite non-IYB group can have only IYB Sylow subgroups. These are statement-only questions requiring human mathematical review.
\subsection{Sources}
\begin{itemize}
\item Local registry: \texttt{data/unsolvedmath/problems.json}, permanent UnsolvedMath ID \(30004309\), title \emph{Solvability and Realizability of Finite Skew Braces}, source problem code \texttt{catalogue entry}, in the Oberwolfach Reports Open Problems collection. No individual proposer is supplied.
\item Supplied report citation: \emph{Mini-Workshop: Algebraic Tools for Solving the Yang--Baxter Equation (2020)}, Oberwolfach Reports, pp.~3207--3242, DOI \texttt{10.4171/owr/2019/51}. \url{https://doi.org/10.4171/owr/2019/51}. The record gives proposed year 2019, whereas the report title supplied in the citation says 2020; these dates are retained without resolving their relationship. The background describes the collection as drawn from workshop reports and their open-problem discussions.
\item Damele, \emph{On the multiplicative group of a two-sided skew brace}. \url{https://arxiv.org/abs/2603.24637}.
\item Cedo and Okniński, \emph{New classes of IYB groups}. \url{https://arxiv.org/abs/2502.14805}.
\item Vendramin, \emph{Problems on skew left braces}. \url{https://www.advgrouptheory.com/journal/Volumes/7/Vendramin%20-%20Problems%20on%20skew%20left%20braces.pdf}.
\item Literature assessment checked at \texttt{2026-08-22}: this exact date and the assessment are supplied dated curation, not independently verified. The record labels the questions open and states that recent work proves substantial solvable and IYB subclasses while the general multiplicative-group, class-two two-sided-brace, and non-IYB Sylow questions remain unresolved. The background characterizes this as dated literature triage derived from the supplied OWR curation and requests independent verification of time-sensitive status and cited claims.
\item The registry labels the statement \texttt{corrected\_verified} and reports manual semantic, typographic, and LaTeX review for self-containment and preservation of the mathematical target. This is supplied review metadata, not independent verification of the source, mathematical correctness, or current problem status.
\end{itemize}
\end{document}
